\printedpage{39}
% head: A. Grothendieck : Dérivateurs, Ch. XII, Édition des Universités de Montpellier II et Paris VII

\msp{page 53}

\section*{IV\quad On axiom W(4) and its variants.}

(\footnote[60]{Cf. also part VI (HOT$_W$), p. 186-196}). I have spent days trying to prove that all these variants are equivalent. I am going to give a list of those, among these variants, which in retrospect still seem to me to be of some use. (With notations different from those on pages 13-14.)

\begin{itemize}[leftmargin=5em,labelsep=0.5em]
\item[\fbox{W(4a)}] Every $W$-fibration (cf. th. 3, p. 51, def. 3, p. 52) is ultrastrong, which amounts to this: for
\[
\begin{tikzcd}
X & X' \arrow[l,"q"'] \arrow[d,"f'"] \\
Y \arrow[from=u,"f"'] & Y' \arrow[l,"p"']
\end{tikzcd}
\]
cartesian in Cat, if $f \in \underline{\Fib}_W$, $p \in W$, then $q \in W$.

\item[W(4a$'$)] As in a), assuming $f \in \Parf_\omega$, i.e. $f$ $W_\omega$-perfect (hence a $W_\omega$-fibration), where $W_\omega \subseteq W$ is the smallest subset $\Sigma$ of Fl(Cat) satisfying W(1, 2 bis, 3). % [sic] petits -> petite, accord avec « partie »

\item[\fbox{W(4b)}] Consider a commutative diagram
$\begin{tikzcd}[ampersand replacement=\&, column sep=small, row sep=small, baseline=(X.base)]
|[alias=X]| X \arrow[rr,"f"] \arrow[dr,"p"'] \& \& Y \arrow[dl,"q"] \\
\& S \&
\end{tikzcd}$
in Cat, with $p, q \in \Fib_W$, $f \in W$.
Then $\forall s \in S$, $f_s : X_s \longrightarrow Y_s$ \texttt{[is in} $W$\texttt{]}.

\item[W(4b$'$)] As in b), assuming \texttt{[in addition]} $X$ and $Y$ split fibered over $S$, $f$ compatible with \texttt{[the]} splittings (hence a cartesian functor).

\item[W(4b$''$)] As in b), assuming $p$, $q$ $W$\textit{-perfect.}

\item[\fbox{W(4c)}] Let $f : X \longrightarrow Y$, with $f \in \Fib_W \cap W$, then the fibers of $f$ are $W$-aspheric.

\item[\dashedbox{W(4c$'$)}] As in c), assuming $f \in \Parf_\omega \cap W$. (This is the old W(4b$''$), stripped of the $W_\infty$-parasite.)
\end{itemize}

\msp{page 54}

I have circled the notions \texttt{[boxed here]} that have been really useful. In practice, one can bet that it is the strongest notion, W(4a), which is the only important one. Here are the tautological implications (with the exception of W(4a$'$) $\Longrightarrow$ W(4b), which is proved in the proof of proposition 4 (pages 6-10) (\footnote[61]{This proof uses a stronger notion of $W$-fibration than that of theorem 3, p. 51.})).

\printedpage{40}
% head: A. Grothendieck : Dérivateurs, Ch. XII, Édition des Universités de Montpellier II et Paris VII

\[
\begin{tikzcd}[column sep=small]
& \mathrm{a} \arrow[d,Rightarrow] & \\
& \mathrm{a}' \arrow[dd,Rightarrow,"{\text{\rule[0.6ex]{2cm}{0.4pt}\ \begin{tabular}[t]{@{}l@{}}doubtful implication,\\ one must weaken\\ b) a little, by\\ strengthening the\\ notion of\\ $W$-fibration\\ used\\ to state b)\end{tabular}}}"] & \\
& & \\
& \mathrm{b} \arrow[dl,Rightarrow] \arrow[dr,Rightarrow] \arrow[dd,Rightarrow] & \\
\mathrm{b}' & & \mathrm{b}'' \\
& \mathrm{c} \arrow[d,Rightarrow] & \\
& \mathrm{c}' &
\end{tikzcd}
\]

I have mainly sweated over showing that c$'$) $\Longrightarrow$ c), or c) $\Longrightarrow$ b), without getting there. The effort has allowed me to get rid of various false ideas, and to replace them with right ones, cf. below. There is all the same a positive result that stays afloat:

\textbf{Proposition 1.} \textit{Condition b) (and probably even b$'$) and b$''$)} (\footnote[62]{and assuming in the case of b$''$, that $W$ satisfies the condition on direct limits W(6 bis), cf. prop. 5 of the ``corollary'' of the, page 14.})\textit{, but I am not going to prove it here \dots) probably implies a) for $\bar{W}$, the strong saturation of $W$. In any case, if $W$ is strongly saturated, conditions a), a$'$), b) (and surely also b$'$) and b$''$)) are equivalent.}

\textsc{Proof} of b) $\Longrightarrow$ a), when $W$ is strongly saturated. In fact, assuming b) (and without even assuming $W$ strongly saturated), we are going to prove that if we have
\[
\begin{tikzcd}
X & X' \arrow[l,"q"'] \arrow[d,"f'"] \\
Y \arrow[from=u,"f"'] & Y' \arrow[l,"p"']
\end{tikzcd}
\]
cartesian in Cat, with
\msp{page 55}
$f$ a $W$-fibration, and $p \in W$, then $\underline{q \in \bar{W}}$, i.e. $\hot_W(q)$ \texttt{[is an]} isomorphism. For this, we interpret the arrow $q$ in $\Hot_W$ in the following way. Let
\[
\xi = \hot_{Y,W}(X/Y), \qquad \xi' = p^*(\xi),
\]
\printedpage{41}
% head: A. Grothendieck : Dérivateurs, Ch. XII, Édition des Universités de Montpellier II et Paris VII
which says that $q$ is identified with the canonical arrow
\[
H_\bullet^{\HOT_W}(Y',\xi') \longrightarrow H_\bullet^{\HOT_W}(Y,\xi) \qquad \text{in } \Hot_W = \HOT_W(e), \tag*{$(*)$}
\]
where we note moreover that $\xi \in \HOT_W(Y)$ is a coefficient \textit{locally constant} on $Y$.

Let us first establish this interpretation of $q$. Note that since $X \in \underline{\Fib}_W Y$, the inverse image $\xi' = p^*(\xi)$ is none other than $\hot_{Y',W}(X' = X \times_Y Y')$. On the other hand, if $g : Y \longrightarrow e$, $g' : Y' \longrightarrow e$ are the structural arrows, the two sides of $(*)$ are interpreted resp. as $g'_!(\xi')$ and $g_!(\xi)$ (the $g'_!$ and $g_!$ in the sense of the prederivator $\HOT_W$). Now the calculation of $g'_!$, $g_!$ gives us that the two sides are identified resp. with
\[
\hot_W(X'), \qquad \hot_W(X),
\]
and we admit that the arrow
\[
\hot_W(X') \longrightarrow \hot_W(X)
\]
is none other than $\hot_W(q)$. Besides, since $X$ is $W$-fibered over $S$, practically by definition, $\xi$ is in $\HOT^{\mathrm{lc}}_W(S)$. So it will suffice for us to prove the

\msp{page 56}

\textbf{Corollary.} \textit{Suppose that $W$ satisfies W(4b), and let $p : S' \longrightarrow S$ be in $W$, $\xi \in \HOT^{\mathrm{lc}}_W(S)$, $\xi' = p^*(\xi)$. Then}
\[
H_\bullet(S',\xi') \xrightarrow{\ \sim\ } H_\bullet(S,\xi)
\]
\textit{is an isomorphism in} $\Hot_W$.

\textsc{Proof}. It suffices to prove that for every $\eta \in \mathrm{Ob}\,\Hot_W$ we have a bijection for the $\Hom(-,\eta)$. By the adjunction formula on $S$, $S'$, this map between the Hom in $\Hot_W$ is identified \texttt{[with]}
\[
\Hom_{\HOT_W(S)}(\xi,\eta_S) \longrightarrow \Hom_{\HOT_W(S')}(\xi',\eta_{S'}),
\]
where $\eta_S = p^*(\eta)$, $\eta_{S'} = p'^*(\eta)$ \texttt{[rather} $\eta_S = g^*(\eta)$\texttt{,} $\eta_{S'} = g'^*(\eta)$\texttt{, where} $g : S \longrightarrow e$\texttt{,} $g' : S' \longrightarrow e$ \texttt{are the structural arrows]}, hence $\eta_{S'} = p^*(\eta_S)$, and the arrow $(*)$ is the arrow on the Hom induced by the functor
\[
p^*_{\mathrm{lc}} : \HOT^{\mathrm{lc}}_W(S) \longrightarrow \HOT^{\mathrm{lc}}_W(S').
\]
Now assuming axiom W(4b) (which we have not used so far), the said functor is an \textit{equivalence of categories} if $p \in W$ (cf. VI, p. 204).

\hfill q.e.d.

For fear of forgetting them, I am now going to make a list of the little questions I have elucidated along the way while scribbling.

\begin{itemize}[leftmargin=5em,labelsep=0.5em]
\item[1°)] If $X$, $Y$ are $W$-smooth over $S$, and $f : X \longrightarrow Y$ an $S$-morphism such that $f_s : X_s \longrightarrow Y_s$ is $W$-smooth for every $s \in \mathrm{Ob}\,S$, it does \textit{not} follow that $f$ is $W$-smooth. But this is nevertheless true when $X$ and $Y$ are split fibered over $S$ (which
\printedpage{42}
% head: A. Grothendieck : Dérivateurs, Ch. XII, Édition des Universités de Montpellier II et Paris VII
\msp{page 57}
already implies their $W$-smoothness over $S$), and $f$ compatible with the splittings and a fortiori cartesian. Dual statement for split cofibered categories and the strictly cartesian functors between them that are $W$-proper fiber by fiber: they are $W$-proper.

\item[2°)] I was interested, in a case of split fibered categories $X$, $Y$ over $S$ and a strictly cartesian functor $f : X \longrightarrow Y$, assumed $W$-perfect fiber by fiber, in concluding that $f$ should also be $W$-perfect. (It was a case where in addition one knows that $X$, $Y$ are also $W$-fibered over $S$, i.e. the base change functors $X_{s'} \longrightarrow X_s$, $Y_{s'} \longrightarrow Y_s$ are in $W$. But that was not much use \dots) The question that remained, in 1°), was that of the $W$\textit{-properness} of $f$. But this is hopeless (even for $Y = S$), by virtue of the brilliant % [sic] fonteurs -> foncteurs, lettre omise

\textbf{Proposition 2.} \textit{Suppose that} $S = 0$
$\begin{tikzcd}[ampersand replacement=\&, row sep=tiny, column sep=small, baseline=(Z.base)]
\& a \\
|[alias=Z]| \phantom{0} \arrow[ur] \arrow[dr] \& \\
\& b
\end{tikzcd}$
\textit{, and let $Z$ be a category over $S$. Then $\Phi(Z)$ is proper over $S$ if and only if $Z = \emptyset$.}

Let us then take in 2°) above $X = \Phi(Z)$, $Y = S$, $f : \Phi(Y) \longrightarrow S$ \texttt{[rather} $f : \Phi(Z) \longrightarrow S$\texttt{]} the structural morphism (so the $f_s$ are $W$-perfect), then $f$ is proper only if $X = \emptyset$. Yet, if we take $Z$ $W$-fibered over $S$, $\Phi(Z)$ is so too, and we have a counterexample to 2°).
\end{itemize}

But we have all the same convinced ourselves of the following result:

\textbf{Proposition 3.} \textit{Let $X$, $Y$ be split fibered categories over $S$, $f : X \longrightarrow Y$ a strictly cartesian $S$-functor, such that}
\msp{page 58}
\textit{for every $s \in S$, $f_s : X_s \longrightarrow Y_s$ is} $W$-fibrant. \textit{Then,} if $W$ satisfies W(4a), $f$ \textit{is $W$-fibrant.}

(I was hoping to get that $f$ is $W$-fibrant ``for free'', to prove that W(4c) implies W(4b). But I had to assume W(4a) \dots)